% !TEX root = main.tex
% Section 04: Design.
% This chapter contains the lab's interface designs.

\chapter{Design}

\section{Status}

Designs in this chapter are working documents. Each design becomes
binding only when promoted to an ADR and signed off in
\texttt{docs/decisions/}.

\section{Resonant Core IPC contract}

The Core exposes a single authenticated transport on loopback
(Unix socket on Linux/macOS, named pipe on Windows). The transport
is language-agnostic; clients speak it from any language that can
serialize JSON or CBOR.

\subsection{Message envelope}

\begin{lstlisting}[language=rccosjson,caption={Request envelope.}]
{
  "id": "<uuidv7>",
  "ts": "<rfc3339>",
  "type": "request",
  "method": "capability.execute",
  "auth": {
    "client_id": "<client id>",
    "signature": "<ed25519 over the rest of the message>"
  },
  "body": { ... }
}
\end{lstlisting}

\begin{lstlisting}[language=rccosjson,caption={Response envelope.}]
{
  "id": "<matches request id>",
  "ts": "<rfc3339>",
  "type": "response",
  "ok": true,
  "body": { ... }
}
\end{lstlisting}

The signature is verified by the Core against a per-client
public key. Clients receive a key on first pairing; the private
key never leaves the client process.

\subsection{Capability request}

\begin{lstlisting}[language=rccosjson,caption={Capability execute body.}]
{
  "capability": "filesystem.read",
  "args": { "path": "/Users/you/Documents/report.pdf" },
  "rationale": "user requested: read the report",
  "context_refs": ["memory:budget-2026"]
}
\end{lstlisting}

The Core validates:
\begin{enumerate}
  \item the client is authorized for the capability;
  \item the args match the capability's typed schema;
  \item the rationale field is non-empty and free of the
        human-only verbs (wallet, sign, payment, publish,
        credential);
  \item the requested path is inside an allow-listed root.
\end{enumerate}

If any check fails, the request is rejected and audited. The AI is
not asked to defend or retry.

\section{Cognitive Core interface}

The Cognitive Core is invoked with a structured request and returns
a structured response. Free-form prose is permitted only as a final
explanation field, never as the primary decision.

\begin{lstlisting}[language=rccosjson,caption={Cognitive Core request.}]
{
  "task": "route",
  "input": "Open the budget spreadsheet and summarize it",
  "available_capabilities": ["filesystem.read", "spreadsheet.inspect"],
  "available_models": ["core-local", "local-7b", "cloud-gpt"],
  "context_summary": "..."
}
\end{lstlisting}

\begin{lstlisting}[language=rccosjson,caption={Cognitive Core response.}]
{
  "intent": "summarize_document",
  "capabilities": ["filesystem.read", "spreadsheet.inspect"],
  "model_tier": "local",
  "memory_queries": ["budget-2026", "spreadsheet:budget"],
  "explanation": "User asked for a summary; the local model can read the file and call the spreadsheet inspect tool."
}
\end{lstlisting}

The Core wraps this in an internal capability request before
executing anything.

\section{Hardware profile}

The hardware profile is a JSON object:

\begin{lstlisting}[language=rccosjson,caption={Hardware profile.}]
{
  "arch": "aarch64",
  "os": "macos",
  "os_version": "15.4",
  "cpu": { "model": "Apple M3", "cores": 8, "features": ["neon"] },
  "ram_mb": 16384,
  "disk_mb": 512000,
  "accelerators": [{ "kind": "metal", "model": "Apple M3 GPU" }]
}
\end{lstlisting}

The model-tier selector is a deterministic function of this object
plus the Cognitive Core's measured properties:

\begin{lstlisting}[language=rccossh,caption={Model tier selector (pseudocode).}]
def select_tier(profile, core_stats):
    if profile.ram_mb < 4096:
        return "ternary-only"
    if profile.ram_mb < 8192:
        return "ternary+small-local"
    if profile.accelerators:
        return "ternary+accelerated-local"
    return "ternary+full-local"
\end{lstlisting}

\section{Capability taxonomy}

Capabilities are namespaced and typed. The first-class namespaces
are:

\begin{itemize}
  \item \texttt{filesystem.*} --- read, write, list, stat.
  \item \texttt{process.*} --- spawn, kill, signal.
  \item \texttt{network.*} --- connect, listen (loopback only by default).
  \item \texttt{browser.*} --- tab, navigation, content (when the
        browser client is present).
  \item \texttt{memory.*} --- query, store, promote.
  \item \texttt{cognition.*} --- route, plan, summarize.
  \item \texttt{device.*} --- notifications, clipboard, audio.
\end{itemize}

Each capability has a JSON schema for its arguments. The Core
rejects requests that do not match the schema.
